\section*{अध्याय \ref{ch:b1:p-block-13-15} --- p-ब्लॉक I: बोरॉन, कार्बन और नाइट्रोजन समूह}

\begin{solution}{exo:b1:p-block-13-15:1}
\ce{NH4NO3}: \ce{NH4+} में $-$III, \ce{NO3-} में $+$V। \ce{N2O5}: $+$V।
\ce{HNO2}: $+$III। \ce{N2H4}: $-$II। \ce{Mg3N2}: $-$III।
\end{solution}

\begin{solution}{exo:b1:p-block-13-15:2}
\ce{BF3}: तीन आबंधों और छह इलेक्ट्रॉनों वाला बोरॉन, त्रिकोणीय समतलीय।
\ce{NH3}: तीन आबंध और एक \omterm{def:b1:lewis-resonance-vsepr:lewis}{एकाकी युग्म}। N का \omterm{def:b1:lewis-resonance-vsepr:lewis}{एकाकी युग्म} B से एक \omterm{def:b1:lewis-resonance-vsepr:lewis-acid}{उपसहसंयोजी
आबंध} बनाता है: \ce{BF3} \omterm{def:b1:lewis-resonance-vsepr:lewis-acid}{लुइस अम्ल} है, \ce{NH3} \omterm{def:b1:lewis-resonance-vsepr:lewis-acid}{लुइस क्षारक};
योगज में दोनों परमाणु चतुष्फलकीय हैं।
\end{solution}

\begin{solution}{exo:b1:p-block-13-15:3}
\ce{CO2 + 2NaOH -> Na2CO3 + H2O}; \ce{SiO2 + 2NaOH -> Na2SiO3 + H2O};
\ce{P4O10 + 12NaOH -> 4Na3PO4 + 6H2O}; \ce{Al2O3 + 6HCl -> 2AlCl3 + 3H2O};
\ce{Al2O3 + 2NaOH + 3H2O -> 2NaAl(OH)4}।
\end{solution}

\begin{solution}{exo:b1:p-block-13-15:4}
$\Delta_r G^\circ = 2 \times (-16.45) = \qty{-32.9}{kJ/mol}$, $\log K =
32.9/5.708 = 5.76$, $K = 10^{5.76}$।
\end{solution}

\begin{solution}{exo:b1:p-block-13-15:5}
छह \ce{SiO3^2-} इकाइयों का वलय: \ce{Si6O18^12-}। द्वि-शृंखला, प्रति चार
सिलिकॉन: दो, दो साझा कोनों के साथ (3 O, प्रत्येक पर आवेश $-2$), दो
तीन के साथ ($2.5$ O, प्रत्येक पर आवेश $-1$): \ce{Si4O11^6-}।
\end{solution}

\begin{solution}{exo:b1:p-block-13-15:6}
\ce{Al(OH)3} का \qty{550}{kg}, $550/78.0 = \qty{7.05}{kmol}$ है, जो
\ce{Al2O3} के \qty{3.53}{kmol}, \qty{360}{kg}, देता है; \qty{92}{\percent} पर:
\qty{331}{kg}।
\end{solution}

\begin{solution}{exo:b1:p-block-13-15:7}
बोरॉन के पास एक रिक्त कक्षक है और वह जल से लिए गए हाइड्रॉक्साइड आयन का
\omterm{def:b1:lewis-resonance-vsepr:lewis}{एकाकी युग्म} ग्रहण करता है: \ce{B(OH)3 + 2H2O <=> B(OH)4- + H3O+}। मुक्त
प्रोटॉन जल के एक अणु से आता है, बोरिक अम्ल से नहीं।
\end{solution}

\begin{solution}{exo:b1:p-block-13-15:8}
N, $-$III से $+$II पर जाता है, प्रति \ce{NH3} पाँच इलेक्ट्रॉन; हर \ce{O2}
चार लेता है। $4 \times 5 = 5 \times 4$: \ce{4NH3 + 5O2 -> 4NO + 6H2O}।
\end{solution}

\begin{solution}{exo:b1:p-block-13-15:9}
\ce{NH3 + HNO3 -> NH4NO3}। एक टन $10^6/80.0 = \qty{1.25e4}{mol}$ है;
उदासीन करने के लिए \qty{1.25e4}{mol} अमोनिया और नाइट्रिक अम्ल बनाने के लिए
\qty{1.25e4}{mol}: $2.50 \times 10^4 \times 17.0 = \qty{425}{kg}$।
\end{solution}

\begin{solution}{exo:b1:p-block-13-15:10}
द्वितीय-आवर्त के परमाणु छोटे होते हैं: उनके \textit{p} कक्षक पार्श्व रूप से अच्छी तरह
अतिव्यापित होकर प्रबल $\pi$ आबंध बनाते हैं, इसलिए \ce{N#N} और \ce{O=C=O} छोटे
स्थायी अणु हैं। तृतीय-आवर्त के परमाणु बड़े होते हैं; उनका $\pi$ अतिव्यापन
कमज़ोर होता है, और वे अधिक एकल आबंध पसंद करते हैं: \ce{P4} चतुष्फलक,
\ce{SiO4} जाल।
\end{solution}

\begin{solution}{exo:b1:p-block-13-15:11}
$E^\circ(\ce{PbO2}/\ce{Pb^2+}) = 1.46 > E^\circ(\ce{Cl2}/\ce{Cl-}) = 1.36$~V:
\ce{PbO2 + 4H+ + 2Cl- -> Pb^2+ + Cl2 + 2H2O}। लेड(IV), लेड(II) की तुलना में
अस्थायी है (अक्रिय युग्म); सिलिकॉन(II) कोई स्थायी अवस्था नहीं है और
\ce{SiO2} \omterm{def:b1:nernst:couple}{ऑक्सीकारक} नहीं है।
\end{solution}

\begin{solution}{exo:b1:p-block-13-15:12}
$160 \times 17.0/14.0 = \qty{194}{Mt}$ अमोनिया, जिसमें $194 \times
3.0/17.0 = \qty{34}{Mt}$ हाइड्रोजन है, जो मुख्यतः प्राकृतिक गैस
(मेथेन और भाप) से बनाया जाता है।
\end{solution}

\begin{solution}{pb:b1:p-block-13-15:1}
\textbf{1.} :N\,$\equiv$\,N: --- एक त्रि-आबंध, अत्यंत प्रबल, बिना किसी ध्रुवता के।
\textbf{2.} 0, $-$III, $+$II, $+$IV, $+$V, और \ce{NH4NO3} में $-$III/$+$V।
\textbf{3.} फलीदार पौधों की जड़ों पर जीवाणु, तड़ित, हेबर--बॉश
प्रक्रम।
\textbf{4.} त्रि-आबंध तोड़ने के लिए ऐसा \omterm{def:b1:elementary-steps:catalyst}{उत्प्रेरक} चाहिए जो पौधों के पास
नहीं है; नाइट्रोजिनेस एंज़ाइम वाले जीवाणुओं के पास है।
\textbf{5.} नाइट्रोजन अपचित होता है (0 से $-$III), हाइड्रोजन ऑक्सीकृत (0 से $+$I)।
\textbf{6.} \ce{N2 + 3H2 <=> 2NH3}।
\textbf{7.} $K = 10^{2 \times 16.45/5.708} = 10^{5.76}$।
\textbf{8.} \qty{25}{\celsius} पर अभिक्रिया बहुत अधिक धीमी है;
\omterm{def:b1:elementary-steps:catalyst}{उत्प्रेरक} केवल गरम होने पर काम करता है, और परिस्थितियाँ एक समझौता हैं जिसकी चर्चा
द्वितीय वर्ष के खंड में है।
\textbf{9.} \qty{58.8}{kmol} अमोनिया: \ce{N2} के \qty{29.41}{kmol},
\qty{824}{kg}; \ce{H2} के \qty{88.2}{kmol}, \qty{176}{kg}।
\textbf{10.} $V(\ce{N2}) = 29.4 \times 10^3 \times 8.314 \times 298.15/10^5 =
\qty{729}{m^3}$; वायु $729/0.78 = \qty{935}{m^3}$।
\textbf{11.} प्राकृतिक गैस से (मेथेन का भाप सुधारण), \ce{CO2}
उपोत्पाद के साथ।
\textbf{12.} मिट्टी में अंतःक्षेपित करने पर वह एक सस्ता उर्वरक है, पर वह एक विषैली
गैस है; ठोस (यूरिया, अमोनियम नाइट्रेट) भंडारण और छिड़काव में अधिक सुगम हैं।
\textbf{13.} N: $-$III $\to$ $+$II (5 इलेक्ट्रॉन), O: 0 $\to$ $-$II (प्रति
\ce{O2} 4): \ce{4NH3 + 5O2 -> 4NO + 6H2O}।
\textbf{14.} \ce{2NO + O2 -> 2NO2}।
\textbf{15.} \ce{3NO2 + H2O -> 2HNO3 + NO}: एक \omterm{def:b1:e-ph-diagrams:disproportionation}{असमानुपातन} (N, $+$IV से
$+$V और $+$II)।
\textbf{16.} \ce{NH3 + 2O2 -> HNO3 + H2O}।
\textbf{17.} $2 \times 58.8 \times 32.0 = \qty{3.76}{t}$।
\textbf{18.} वह ऑक्सीकरण को तीव्र और NO के लिए वरणात्मक बनाता है, अधिक
स्थायी \ce{N2} के बजाय।
\textbf{19.} \ce{NH3 + HNO3 -> NH4NO3}।
\textbf{20.} आधा।
\textbf{21.} \ce{NH4NO3} के \qty{29.4}{kmol}, \qty{2.35}{t}।
\textbf{22.} $28.0/80.0 = \qty{35}{\percent}$।
\textbf{23.} उसमें एक ही \omterm{def:b1:crystals-metals:crystal}{क्रिस्टल} में एक \omterm{def:b1:nernst:couple}{ऑक्सीकारक} (नाइट्रेट) और एक \omterm{def:b1:nernst:couple}{अपचायक} (अमोनियम)
है: गरम करने पर या ईंधनों के साथ मिलाने पर वह विस्फोटक रूप से अपघटित हो सकता है।
\textbf{24.} $58.8 \times 63.0 =$ \textbf{प्रति टन अमोनिया \qty{3.7}{t} नाइट्रिक
अम्ल}।
\end{solution}
